\section{第6章\quad 重抽样方法（Bootstrap 与 Jackknife）}

\subsection{总览}

第 3 章评估估计量性质时假定\textbf{总体分布已知}（能反复从总体抽样）。
现实中往往只有一份样本、总体分布 $F$ 未知。
\textbf{Bootstrap（自助法，Efron, 1979）}的核心创意：
\textbf{用经验分布 $F_n$（即样本本身）代替总体 $F$，从样本中有放回再抽样}，
从而近似估计量 $\hat\theta(X_1,\dots,X_n)$ 的分布及其特征（标准误、偏差、置信区间）。

\subsection{从蒙特卡洛到 Bootstrap}

\begin{kbox}{对比记忆}
\textbf{若 $F(x;\theta)$ 已知}（参数 MC 法）：
\begin{enumerate}[leftmargin=2em]
  \item 从总体 $F$ 中独立生成 $n$ 个数据，算出 $\hat\theta_1$；
  \item 重复 $m$ 次得 $\hat\theta_1,\dots,\hat\theta_m$；
  \item 用 $\bar{\hat\theta}=\frac1m\sum\hat\theta_i$ 和样本方差估计 $\E\hat\theta,\ \Var\hat\theta$。
\end{enumerate}
\textbf{若 $F$ 未知}（Bootstrap）：把"从 $F$ 抽样"替换为"\textbf{从原样本 $\{X_1,\dots,X_n\}$ 中有放回抽 $n$ 个}"：
\begin{enumerate}[leftmargin=2em]
  \item 有放回抽得自助样本 $X^*_{11},\dots,X^*_{1n}$，算 $\hat\theta^*_1$；
  \item 重复 $m$ 次得 $\hat\theta^*_1,\dots,\hat\theta^*_m$；
  \item 用 $\{\hat\theta^*_i\}$ 的分布近似 $\hat\theta$ 的分布。
\end{enumerate}
\end{kbox}

\subsection{Bootstrap 估计标准误与偏差}

\textbf{标准误的 Bootstrap 估计}：
\[
\widehat{se}_{B}=\sqrt{\frac{1}{m-1}\sum_{i=1}^{m}\big(\hat\theta_i^{*}-\bar{\hat\theta}^{*}\big)^2},
\qquad \bar{\hat\theta}^{*}=\frac1m\sum_{i=1}^m\hat\theta_i^* .
\]

\textbf{偏差的 Bootstrap 估计}：真偏差 $\mathrm{Bias}(\hat\theta)=\E\hat\theta-\theta$。
Bootstrap 世界里"真参数"是 $\hat\theta$（由原样本算出），"估计量"是 $\hat\theta^*$，故
\[
\widehat{\mathrm{Bias}}_B=\bar{\hat\theta}^{*}-\hat\theta .
\]
\textbf{偏差修正估计}：
\[
\tilde\theta=\hat\theta-\widehat{\mathrm{Bias}}_B=2\hat\theta-\bar{\hat\theta}^{*}.
\]

\textbf{经典例子}：$\hat\sigma^2=\frac1n\sum(X_i-\bar X)^2$ 的真实偏差为
$\E\hat\sigma^2-\sigma^2=-\sigma^2/n$（低估）。
用 Bootstrap 估计该偏差：对每个自助样本算 $\hat\sigma^{2*}_i$，
$\widehat{\mathrm{Bias}}=\bar{\hat\sigma}^{2*}-\hat\sigma^2$，应接近 $-\hat\sigma^2/n$。

\subsection{Jackknife（刀切法）}

逐个"剔除一个观测"构造 $n$ 个删一样本：
$\hat\theta_{(i)}$ 为去掉 $X_i$ 后用其余 $n-1$ 个数据算的估计，
$\bar\theta_{(\cdot)}=\frac1n\sum_i\hat\theta_{(i)}$。
\[
\widehat{\mathrm{Bias}}_{J}=(n-1)\big(\bar\theta_{(\cdot)}-\hat\theta\big),\qquad
\widehat{se}_{J}=\sqrt{\frac{n-1}{n}\sum_{i=1}^{n}\big(\hat\theta_{(i)}-\bar\theta_{(\cdot)}\big)^2}.
\]
特点：确定性算法（无需随机抽样）、计算量 $n$ 次；
对光滑统计量（均值、方差等）效果好，对不光滑统计量（如中位数）可能失效；
可视为 Bootstrap 的线性近似。

\subsection{Bootstrap 置信区间}

设自助估计 $\hat\theta^*_1,\dots,\hat\theta^*_m$，从小到大排序，
$\hat\theta^*_{(\alpha/2)},\hat\theta^*_{(1-\alpha/2)}$ 为经验分位数。三种常用区间：
\begin{enumerate}[leftmargin=2em]
  \item \textbf{标准正态（标准误）法}：假定 $\hat\theta$ 近似正态，
  \[
  \big[\hat\theta-u_{1-\alpha/2}\,\widehat{se}_B,\ \hat\theta+u_{1-\alpha/2}\,\widehat{se}_B\big].
  \]
  \item \textbf{百分位数法}：直接取自助分布的分位数
  \[
  \big[\hat\theta^*_{(\alpha/2)},\ \hat\theta^*_{(1-\alpha/2)}\big].
  \]
  \item \textbf{枢轴量（基本）法}：用 $\hat\theta^*-\hat\theta$ 的分布近似 $\hat\theta-\theta$ 的分布，得
  \[
  \big[2\hat\theta-\hat\theta^*_{(1-\alpha/2)},\ 2\hat\theta-\hat\theta^*_{(\alpha/2)}\big].
  \]
\end{enumerate}

\begin{wbox}{易混点}
百分位法区间和枢轴法区间\textbf{端点"翻转"}：枢轴法用 $2\hat\theta$ 减去\textbf{反向}分位数。
偏态分布下两者差异明显；正态近似法只在自助分布近似对称正态时合适。
\end{wbox}

\subsection{本章考点与题型}

\begin{ebox}{高频题型}
\begin{enumerate}[leftmargin=2em]
  \item \textbf{叙述 Bootstrap 基本思想}：经验分布替代总体 + 有放回重抽样（概念题）。
  \item \textbf{写出 Bootstrap 估计标准误/偏差的步骤公式}；给小数据集手算一两步。
  \item \textbf{偏差修正}：$\tilde\theta=2\hat\theta-\bar{\hat\theta}^*$ 的来历。
  \item \textbf{Jackknife 公式}及与 Bootstrap 的比较（确定性 vs 随机、$n$ 次 vs $m$ 次、不光滑统计量失效）。
  \item \textbf{三种 Bootstrap 置信区间}的公式与适用场景辨析。
\end{enumerate}
\end{ebox}
